\subsection{Experimental Results}

We evaluate three global-queue-based scheduling strategies---\texttt{GLOBAL\_QUEUE}, \texttt{BATCH\_PULL}, and \texttt{HYBRID\_QUEUE}---on two workloads using the Yahoo Streaming Benchmark (YSB) with 50 million tuples (4.8\,GB).
All experiments run on a dual-socket server with 2 NUMA nodes (32 cores each, 64 total) and 512\,GB RAM.
For thread counts $\leq 32$, processes are pinned to a single NUMA node via \texttt{numactl --cpunodebind=0 --membind=0}.
Each configuration is executed once; batch sizes are swept over $\{1, 2, 4, 8, 16, 32, 64\}$ for \texttt{BATCH\_PULL} and \texttt{HYBRID\_QUEUE}.

\subsubsection{Stateless Workload (Filter Only)}

The stateless workload applies a simple predicate (\texttt{WHERE event\_type < 1}) with approximately 80\% selectivity.
This workload is compute-light and dominated by buffer throughput, making it sensitive to scheduling overhead and queue contention.

Figure~\ref{fig:heatmap-stateless} shows the throughput heatmap for all three strategies.
The key observations are:

\begin{itemize}
    \item \textbf{HYBRID\_QUEUE achieves the highest throughput overall}, peaking at \textbf{1,028\,MB/s at 16 threads} with batch size 4, corresponding to a \textbf{13.3$\times$ speedup} over a single thread.
    This is because at 16 threads (one full NUMA node), the combination of batched pulls and local successor queues minimizes both shared queue contention and cross-core cache invalidation.

    \item \textbf{Batch size 4--8 is optimal for HYBRID\_QUEUE.}
    Smaller batches (1--2) increase contention on the global queue, while larger batches ($\geq 16$) cause load imbalance as threads hoard work.
    The sweet spot at 4--8 balances contention reduction with fair work distribution.

    \item \textbf{At 32 threads, throughput drops due to hyperthreading contention.}
    NUMA node~0 has 16 physical cores (CPUs 0--15) and 16 hyperthreads (CPUs 32--47).
    At 16T, each thread runs on a dedicated physical core.
    At 32T, threads share physical cores via hyperthreading, competing for execution units and L1/L2 cache.
    The best 32T configuration is HYBRID\_QUEUE BS=8 at 817\,MB/s (10.6$\times$), a 20\% drop from the 16T peak.
    At 64 threads (no NUMA pinning), cross-NUMA memory access adds further overhead, with HYBRID\_QUEUE BS=8 achieving 599\,MB/s (7.8$\times$).

    \item \textbf{BATCH\_PULL performs similarly to GLOBAL\_QUEUE.}
    Since the stateless workload produces no successor tasks (single-pipeline query), the per-thread queue advantage of HYBRID\_QUEUE comes solely from the batched pull reducing contention---the local successor queue is unused.
    Yet HYBRID\_QUEUE still outperforms BATCH\_PULL, suggesting that the per-thread queue structure itself provides better cache behavior for the pull operation.

    \item \textbf{GLOBAL\_QUEUE peaks at 966\,MB/s at 16T}, competitive with BATCH\_PULL (1,025\,MB/s at BS=4) but 6\% below HYBRID\_QUEUE.
    At higher thread counts, GLOBAL\_QUEUE degrades more than the batched alternatives due to increasing atomic contention on the single shared queue.
\end{itemize}

\subsubsection{Stateful Workload (Aggregation with Tumbling Window)}

The stateful workload performs a grouped aggregation with a 30-second tumbling window (\texttt{GROUP BY campaign\_id WINDOW TUMBLING(current\_ms, SIZE 30 SEC)}).
This workload is compute-heavy due to hash table lookups for 10,000 distinct groups and window state management.

Figure~\ref{fig:heatmap-stateful} shows the throughput heatmap.
The key observations are:

\begin{itemize}
    \item \textbf{Scaling is limited to approximately 2.4$\times$}, with the best configuration achieving 205\,MB/s at 32 threads (\texttt{GLOBAL\_QUEUE}).
    The aggregation operator's shared window state and hash table become the bottleneck, not the scheduling strategy.

    \item \textbf{All three strategies perform within 10--15\% of each other} across all thread counts and batch sizes.
    This confirms that for compute-bound stateful queries, scheduling overhead is negligible compared to operator execution time.

    \item \textbf{GLOBAL\_QUEUE wins at 32T} (205\,MB/s), likely because the stateful pipeline produces very few successor tasks (one window emission per 30 seconds of data, i.e., approximately 1,667 emissions across 50 million tuples).
    Without significant successor task volume, HYBRID\_QUEUE's local queue advantage is unused.

    \item \textbf{Batch size has minimal impact} on the stateful workload.
    Throughput varies by less than 10\% across all batch sizes for both BATCH\_PULL and HYBRID\_QUEUE, confirming that queue contention is not the bottleneck.

    \item \textbf{The throughput dip at 16T} (152--154\,MB/s, below the 8T result of 155--162\,MB/s) is likely caused by increased contention on the shared aggregation hash table: with 16 threads all updating the same window state concurrently, synchronization overhead exceeds the parallelism gain.
    At 32T, throughput recovers (205\,MB/s) because the larger number of threads amortizes the per-thread overhead despite hyperthreading contention.
\end{itemize}

\subsubsection{Summary}

\begin{table}[h]
\centering
\caption{Best configuration per thread count for both workloads.}
\label{tab:best-configs}
\begin{tabular}{r|lr|lr}
\hline
\textbf{Threads} & \multicolumn{2}{c|}{\textbf{Stateless}} & \multicolumn{2}{c}{\textbf{Stateful}} \\
 & MB/s & Config & MB/s & Config \\
\hline
1  & 77   & all similar       & 85   & all similar \\
2  & 152  & HYBRID BS=32      & 130  & BATCH BS=64 \\
4  & 298  & HYBRID BS=32      & 163  & BATCH BS=32 \\
8  & 595  & HYBRID BS=8       & 162  & BATCH BS=32 \\
16 & 1028 & HYBRID BS=4       & 154  & BATCH BS=4 \\
32 & 817  & HYBRID BS=8       & 205  & GLOBAL \\
64 & 599  & HYBRID BS=8       & 166  & BATCH BS=1 \\
\hline
\end{tabular}
\end{table}

The results show that the optimal scheduling strategy depends on the workload characteristics:

\begin{itemize}
    \item For \textbf{stateless, throughput-bound workloads}, \texttt{HYBRID\_QUEUE} with batch size 4--8 provides the best performance, achieving up to 13.3$\times$ speedup.
    The combination of batched global queue pulls and per-thread local queues minimizes contention while preserving cache locality.

    \item For \textbf{stateful, compute-bound workloads}, the scheduling strategy matters less.
    \texttt{GLOBAL\_QUEUE} or \texttt{BATCH\_PULL} with any batch size performs comparably, as the operator execution time dominates.

    \item \textbf{Batch size 4--8 is a robust default} that works well across both workloads and thread counts.
    It provides enough batching to reduce contention without causing significant load imbalance.

    \item \textbf{Hardware topology awareness is critical.}
    The throughput drop from 16T (dedicated physical cores) to 32T (hyperthreading) is approximately 20\%, and the further drop to 64T (cross-NUMA) adds another 25\%.
    This highlights the need for both hyperthreading-aware and NUMA-aware scheduling in future work.
\end{itemize}
